\documentclass{ctexbeamer}

% 使用 <thubeamer> 主题
% 模板选项如下
% (a.1) smoothbars: 页面顶端单行显示目录，默认选项
\usetheme[sectiontoc]{thubeamer}

\usepackage{media9}
\usepackage{attachfile2}
\usepackage{listings}
\usepackage{multirow}
\usepackage{colortbl}
\usepackage{booktabs}
\usepackage{tikz}
\usetikzlibrary{shapes.geometric, arrows.meta, positioning, fit, backgrounds, calc}
\usepackage{animate}
\lstnewenvironment{textcode}{
  \lstset{
    basicstyle=\footnotesize\ttfamily,
    breaklines=true,
    frame=single,
    numbersep=2pt,
    xleftmargin=0.5em,
    framesep=1mm
  }
}{}


% (a.2) sidebar: 页面左侧分栏显示目录
% \usetheme[sidebar]{thubeamer}

% (b) sectiontoc: 在每节（section）前显示目录，并高亮显示当前节，默认不显示

% (c) subsectiontoc: 在每小节（subsection）前显示目录，并高亮显示当前节和当前小节，默认不显示

% (d) en: 仅使用英文来制作 beamer，应用此选项后，汉字将全部无法编译

% 图片存放路径

\graphicspath{{figures/}}

\include{macros}

% 封面信息，方括号内容是显示在左侧边栏的内容（当选择 sidebar 主题时有效）
\title[篮球高光检测]{篮球高光检测项目\\[2mm] 汇报}
\author[唐梓楠]{牛衍昌\\鲁继元\\唐梓楠\\[5mm] 导师：黄必清教授}
\institute[清华大学]{\small 清华大学}
\date{\small \vskip -10pt \today}

% 开始写文章
\begin{document}

% 标题页
\begin{frame}
  \maketitle
\end{frame}

% 目录页

\begin{frame}{实现长时间追踪（>1min）}
  优化方法：采用滑动窗口，先把视频切分成多个窗口，每个窗口内进行追踪，最后再利用重叠的窗口，进行窗口间的匹配。

  最优窗口大小还在调试中。

  \textattachfile{demo.mp4}{demo.mp4}
\end{frame}

\begin{frame}{SigLip2 的替换}
  仅用 SigLip2 / Qwen3-VL-Embedding-2B 的效果不太行，ViT 在细腻度区分上具有局限性。

  \textattachfile{demo2.mp4}{demo2.mp4}
\end{frame}

\begin{frame}{技术栈更新}
  \begin{itemize}
    \item 实现了 Yolo12x 到 Yolo26 的替换
    \item 改为用 uv 进行 Python 项目管理，review 了代码，可以提交到鹰眼的代码仓库
  \end{itemize}
\end{frame}

\begin{frame}{后续工作}
  \begin{itemize}
    \item 优化窗口大小等超参数
    \item 将 ViT 方法与人脸 ReID 方法结合，例如人脸遮挡时，可以利用衣服的相似度进行 ReID，人脸存在时，可以利用人脸的相似度进行 ReID
    \item 优化查询逻辑，例如先用查询图查询起始片段，再利用起始片段的查询结果优化后续片段的查询
  \end{itemize}
\end{frame}

% 结束文档撰写
\end{document}
